\documentclass[a4paper,11pt]{article}
\usepackage[utf8]{inputenc}
\usepackage[T1]{fontenc}
\usepackage[french]{babel}
\usepackage{lmodern}
\usepackage{amsmath,amssymb,amsthm}
\usepackage{mathrsfs}
\usepackage{enumitem}
\usepackage{multicol}

\setlist[enumerate]{label=\textbf{\textcolor{Couleur8}{\arabic*.}}, left=1em}

% Si vous voulez aussi modifier les listes enumerate imbriquées :
\setlist[enumerate,2]{label=\textcolor{Couleur8}{\alph*)}}
\setlist[enumerate,3]{label=\textcolor{Couleur8}{\roman*)}}
\setlist[enumerate,4]{label=\textcolor{Couleur8}{\Alph*)}}
\usepackage{graphicx}
\usepackage{xcolor}
\usepackage{tcolorbox}
\usepackage{geometry}
\usepackage{fancyhdr}
\usepackage{titlesec}
\usepackage{cancel}
\usepackage{ifthen}
\usepackage{tikz}
\usetikzlibrary{arrows.meta}
\usepackage{tkz-tab}
\usepackage{eurosym}

% OPTION PROF/ÉLÈVE
% Décommentez UNE des deux lignes suivantes :
\newboolean{prof}\setboolean{prof}{true}   % Version professeur (avec corrections des devoirs)
\newboolean{prof}\setboolean{prof}{true}  % Version élève (sans corrections des devoirs)

\definecolor{Couleur1}{HTML}{4A5D7A} %Th
\definecolor{Couleur2}{HTML}{A5B5D6} %Prop
\definecolor{Couleur3}{HTML}{A8C68A} %Méthode
\definecolor{Couleur6}{HTML}{8A9CC1} %Def
\definecolor{Couleur7}{HTML}{E6B459} %Rq Voc
\definecolor{Couleur8}{HTML}{D36740} %Titres
\definecolor{correction}{HTML}{D36740} % Couleur pour les corrections
\definecolor{coopmathscorr}{HTML}{F15929} % Couleur corrections coopmaths

% Configuration des marges
\geometry{
	top=2cm,
	bottom=2cm,
	inner=1.5cm,
	outer=1.5cm,
}

% Police
\renewcommand{\familydefault}{\sfdefault}

% Compteurs
\newcounter{seancenum}
\newcounter{seancenumD}
\setcounter{seancenum}{1}
\setcounter{seancenumD}{1}

% Configuration des en-têtes et pieds de page
\pagestyle{fancy}
\fancyhf{}
\ifthenelse{\boolean{prof}}{
    \fancyhead[L]{\textcolor{Couleur8}{\textbf{Corrections Automatismes - Classe de seconde - Version Professeur}}}
}{
    \fancyhead[L]{\textcolor{Couleur8}{\textbf{Corrections Devoirs - Séance 12 - Classe de seconde}}}
}
\fancyhead[R]{\textcolor{Couleur8}{\textbf{2025-2026}}}
\fancyfoot[L]{\textcolor{Couleur8}{Mme Bertail et Mme Pinto}}
\fancyfoot[R]{\textcolor{Couleur8}{\thepage}}
\renewcommand{\headrulewidth}{1pt}
\renewcommand{\headrule}{\hbox to\headwidth{\color{Couleur8}\leaders\hrule height \headrulewidth\hfill}}

% Environnements personnalisés
\newtcolorbox{seance}[1]{
    colback=Couleur6!10,
    colframe=Couleur1!65,
    fonttitle=\bfseries\large,
    title=Correction Séance \theseancenum #1,
    left=5mm,
    right=5mm,
    top=3mm,
    bottom=3mm,
    arc=0mm,
    after=\stepcounter{seancenum}
}

\newtcolorbox{seanceD}[1]{
    colback=Couleur6!5,
    colframe=Couleur6!45,
    fonttitle=\bfseries\large,
    title=Correction Devoirs - Séance \theseancenumD #1,
    left=5mm,
    right=5mm,
    top=3mm,
    bottom=3mm,
    arc=0mm,
    after=\stepcounter{seancenumD}
}

\begin{document}

\setcounter{seancenumD}{12}
\newpage
\begin{seanceD}{} %12
\begin{enumerate}
\item
 $f$ est une fonction affine, elle a donc une expression de la forme  $f(x)=ax+b$ avec $a$ et $b$ des nombres réels. D'après le cours, on sait que pour $u\neq v$, $a=\dfrac{f(u)-f(v)}{u-v}$. Avec $u=9$ et  $v=7$, on obtient  :  $a=\dfrac{f(9)-f(7)}{9-7}=\dfrac{32-26}{9-7}=\dfrac{6}{2}=3$. On en déduit que la fonction $f$ s'écrit sous la forme :    $f(x)=3 x +b.$\\On obtient $b$ en utilisant (au choix) une des deux données de l'énoncé, par exemple $f(9)=32$. Comme $f(x)=3x +b$, alors $f(9)=
              3\times9+b=27+b$ . On en déduit :\\$\begin{aligned}f(9)=32&\iff 27+b=32\\&\iff b=5\\\end{aligned}$\\On en déduit $f(x)=$ ${\color[HTML]{f15929}\boldsymbol{3x+5}}$.


\item \begin{enumerate}
\item On cherche les réels qui sont dans  $]1;27]$ ou bien  $]6;13[$, ou dans les deux. On regarde donc la partie de l'intervalle qui est coloriée, soit en bleu, soit en rouge, soit en bleu et rouge :\\On a $]6;13]\subset ]1;27]$ donc $I={\color[HTML]{f15929}\boldsymbol{]1;27]}}$.\\\begin{tikzpicture}

    \tikzset{
      point/.style={
        thick,
        draw,
        cross out,
        inner sep=0pt,
        minimum width=2pt,
        minimum height=2pt,
      },
    }
    \clip (-2,-0.7) rectangle (15,0.5);
    	\draw[color ={black},line width = 3] (0,0)--(12,0);
	\draw[color ={red},line width = 3] (2,-0.1)--(9,-0.1);
	\draw[color ={blue},line width = 3] (5,0.1)--(7,0.1);
	\draw[very thick,color={red}] (1.8,0.2)--(2,0.2)--(2,-0.2)--(1.8,-0.2);
	\draw[color={red}] (2,-0.2) node[below] {$1$};
	\draw[very thick,color={red}] (8.8,0.2)--(9,0.2)--(9,-0.2)--(8.8,-0.2);
	\draw[color={red}] (9,-0.2) node[below] {$27$};
	\draw[very thick,color={blue}] (4.8,0.2)--(5,0.2)--(5,-0.2)--(4.8,-0.2);
	\draw[color={blue}] (5,-0.2) node[below] {$6$};
	\draw[very thick,color={blue}] (6.8,0.2)--(7,0.2)--(7,-0.2)--(6.8,-0.2);
	\draw[color={blue}] (7,-0.2) node[below] {$13$};

\end{tikzpicture}
\item On cherche les réels qui sont à la fois dans $]-\infty;22]$ et dans $]31;40]$.On regarde la partie de l'intervalle qui est coloriée à la fois en bleu et en rouge. On observe que les deux intervalles sont disjoints donc aucun réel n'appartient aux deux ensembles.\\$I={\color[HTML]{f15929}\boldsymbol{\emptyset}}$.\\\begin{tikzpicture}

    \tikzset{
      point/.style={
        thick,
        draw,
        cross out,
        inner sep=0pt,
        minimum width=2pt,
        minimum height=2pt,
      },
    }
    \clip (-2,-0.7) rectangle (15,0.5);
    	\draw[color ={black},line width = 3] (0,0)--(12,0);
	\draw[color ={red},line width = 3] (0,0)--(5,0);
	\draw[color ={blue},line width = 3] (6,0)--(10,0);
	
	\draw[very thick,color={red}] (4.8,0.2)--(5,0.2)--(5,-0.2)--(4.8,-0.2);
	\draw[color={red}] (5,-0.2) node[below] {$22$};
	\draw[very thick,color={blue}] (5.8,0.2)--(6,0.2)--(6,-0.2)--(5.8,-0.2);
	\draw[color={blue}] (6,-0.2) node[below] {$31$};
	\draw[very thick,color={blue}] (9.8,0.2)--(10,0.2)--(10,-0.2)--(9.8,-0.2);
	\draw[color={blue}] (10,-0.2) node[below] {$40$};

\end{tikzpicture}
\end{enumerate}

\item  \textbf {a.}  Résolution de l'inéquation :\\
              $\begin{aligned}
              6x+3&<x+5\\
              6x+3\,{\color[HTML]{f15929}\boldsymbol{-x}}&<x+5\,{\color[HTML]{f15929}\boldsymbol{-x}}\\
              5x+3&<5\\
              5x+3\,{\color[HTML]{f15929}\boldsymbol{-3}}&<5\,{\color[HTML]{f15929}\boldsymbol{-3}}\\
              5x&<2\\
            x&<\dfrac{2}{5}
              \end{aligned}$
              \\
              L'  ensemble $S$ des solutions de l'inéquation est
              $S= \left]-\infty\,;\,\dfrac{2}{5}\right[$.\\\textbf {b.}  
              La courbe $\mathscr{C}_f$ est en dessous de la courbe $\mathscr{C}_g$ sur l'intervalle $\left]-\infty\,;\,\dfrac{2}{5}\right[$.



\item  $A =  (9k+4)(-7k-3)$ \\$A =  9k\times (-7k) + 
  9k\times (-3) +
  4 \times (-7k)  +
  4 \times (-3)$ \\$A = -63k^2 + 
  (-27k) +
  (-28k) +
  (-12)$ \\$A =$ ${\color[HTML]{f15929}\boldsymbol{ -63k^2-55k-12 }}$


\item \begin{enumerate}
\item On utilise la formule du cours qui exprime le taux d'évolution $t$ en fonction de la valeur initiale $V_i$ et la valeur finale $V_f$ : $t=\dfrac{V_f-V_i}{V_i}=\dfrac{43\,680-52\,000}{52\,000}=-0{,}16=\dfrac{-16}{100}$.\\Le taux d'évolution de la population est ${\color[HTML]{f15929}\boldsymbol{-16}}~\%$

\medskip
Méthode $2$ : On arrive aussi au même résultat en passant par le coefficient multiplicateur : \\$CM=\dfrac{V_f}{V_i}=\dfrac{43\,680}{52\,000} = 0{,}84$ donc $t = 0{,}84-1= -16~\%$
\item Augmenter de $60~\%$ revient à multiplier par $1 + \dfrac{60}{100} = 1+ 0{,}6 = 1{,}6$.\\Pour retrouver le prix initial, on va donc diviser le prix final par $1{,}6$ donc $\dfrac{15{,}84}{1{,}6}  = 9{,}90$\\Avant l'augmentation cet article coûtait ${\color[HTML]{f15929}\boldsymbol{9{,}90}}$\,\euro{}.
\end{enumerate}



\end{enumerate}
\end{seanceD}

\end{document}